\documentclass{article}
\usepackage{ctex}
\usepackage{listings}
\usepackage{xcolor}
\usepackage{geometry}
\usepackage{amsmath}  % 修复数学公式命令
\geometry{a4paper, margin=1in}

% 代码高亮
\lstset{
    language=Python,
    basicstyle=\small,
    commentstyle=\color{gray},
    stringstyle=\color{red},
    keywordstyle=\color{blue},
    frame=single,
    breaklines=true,
    numbers=left,
    numberstyle=\tiny\color{gray},
    showstringspaces=false
}

\title{CTF Crypto 题解：primeRSA}
\author{}
\date{}

\begin{document}
\maketitle

\section{题目信息}
\subsection{题目源码}
题目给出加密代码如下，使用 RSA 加密，模数 $n$ 由 \texttt{getPrime(1024)} 生成，公钥指数 $e=0$x$10001=65537$。

\begin{lstlisting}
from Crypto.Util.number import getPrime, bytes_to_long
import re

flag = open("flag.txt", "rb").read().strip()
m = bytes_to_long(flag)
assert re.fullmatch(r"flag{[A-Za-z0-9_]+}", flag.decode())
n = getPrime(1024)
e = 0x10001
c = pow(m, e, n)

print(f"{n=}")
print(f"{e=}")
print(f"{c=}")
\end{lstlisting}

程序运行后输出三行内容至 \texttt{output.txt}：
\begin{itemize}
    \item 第一行：\texttt{n=大整数}
    \item 第二行：\texttt{e=65537}
    \item 第三行：\texttt{c=密文大整数}
\end{itemize}

\section{原理分析}
标准 RSA 算法中，若模数 $n$ 为两个大素数乘积 $n=p\cdot q$，则欧拉函数：
$$
\varphi(n) = (p-1)(q-1)
$$
私钥 $d$ 满足：
$$
e\cdot d \equiv 1 \pmod{\varphi(n)}
$$

本题漏洞点：
题目中直接使用 \texttt{getPrime(1024)} 生成 \textbf{单个素数} 作为模数 $n$，即 $n$ 本身是素数。
此时欧拉函数简化为：
$$
\varphi(n) = n-1
$$

解密公式：
\[
m = c^d \pmod{n},\quad d = e^{-1} \pmod{n-1}
\]
将解密得到的大数 $m$ 转为字节串，即可得到 \texttt{flag}。

\section{解题过程}
\subsection{解密脚本}
使用 SageMath 运行，自动读取 \texttt{output.txt} 并完成解密、反向验证，脚本如下：
\begin{lstlisting}
# solve.py
# 运行命令：sage solve.py
with open("output.txt", "r") as f:
    lines = f.read().splitlines()

n = int(lines[0].split("=")[1])
e = int(lines[1].split("=")[1])
c = int(lines[2].split("=")[1])

phi = n - 1
d = pow(e, -1, phi)
m = pow(c, d, n)

flag = bytes.fromhex(hex(m)[2:])

m_check = int.from_bytes(flag, "big")
c_check = pow(m_check, e, n)

print("解密得到 flag =", flag)
print("原密文 c     =", c)
print("验证密文 c   =", c_check)
print("验证结果     =", c == c_check)
\end{lstlisting}

\subsection{运行环境与命令}
运行环境：WSL Ubuntu + SageMath
执行命令：
\begin{lstlisting}[language=bash]
sage solve.py
\end{lstlisting}

\end{document}